\section{Limitations and Ethical Considerations}
\label{sec:limitations}

\paragraph{Domain.} The corpus is synthetic: every training and evaluation track was generated by MiniMax Music 3 from a narrow set of genre captions, and generalisation to recorded music is not established. Anecdotally the encoder is usable on real references, but we report no metric for it.

\paragraph{Open validation.} Three measurements would bound the evidence and were not available to us, because they require the official language model, depth decoder and condition encoder as a scoring path. First, a floor for condition-replay cosine: random codes, codes from another track, and tuples with only the semantic or only the acoustic codes correct would calibrate what 0.87 means against 0.66 for the community encoder and 0.9999 for the true codes. Second, a study on recorded music, scoring functional preservation of a reference (embedding, chroma, tempo or listening judgements) after generation, since no ground-truth codes exist for real audio. Third, an end-to-end check that higher replay cosine corresponds to better preservation after the diffusion and decoding stages. Two matched ablations would sharpen the architectural claim: a 512-wide model trained with the v2 schedule, and a parameter-matched independent-head or non-causal decoder against v4.

\paragraph{Context and decoding.} The context is 128 frames (5.12 s) with no state across windows. The depth decoder has a large exposure gap between teacher-forced and greedy decoding, and we tried no beam search, scheduled sampling~\cite{bengio2015scheduled} or differentiable depth sampling. Each architecture was trained once.

\paragraph{Metrics and teacher.} Condition-replay cosine stops before the diffusion transformer and is not an audio-quality measure; exact-token agreement does not identify functionally equivalent code tuples; and the teacher distributions come from the generator's text-conditioned prior, not from an audio-conditioned posterior.

\paragraph{Ethical considerations.} No original encoder weights or source were used; the encoder is trained only on outputs of the released model, obtained through its public interface, and it does not reproduce any withheld weights. It enables reference-conditioned generation, which can be used to imitate a recording's structure; the released integration constrains only the semantic stream and leaves timbre to the generator, and use of the weights and of any generated audio remains subject to the MiniMax Music 3 terms. The corpus contains no recorded music. The v3 run's provenance additionally falls under the MERT license.
